\chapterbackmatter{Solutions to Supplementary Exercises}

\begin{enumerate}
\SupplementarySolution{1}{A Census of One-Loop QED Graphs}
The tree graph contains one virtual photon and two fermion--photon vertices.
At order \(e^4\), the connected loop topologies are:
\begin{enumerate}
\item a photon vacuum-polarization insertion (with every charged flavor that
can circulate in the loop);
\item a one-loop vertex correction at the electron or muon annihilation
vertex;
\item self-energy insertions on each of the four external fermion legs; and
\item the box and crossed-box two-photon-exchange graphs.
\end{enumerate}
There are also counterterm graphs for the photon propagator, both vertices,
and the fermion two-point functions.  In an on-shell scheme the external-leg
self-energies combine with their counterterms and LSZ factors to leave unit
pole residue.  The vacuum-polarization and vertex subgraphs are replaced by
their subtracted, finite versions, while the box graphs are ultraviolet
finite after their orderings are summed.  Thus a convenient organization is
\[
 \mathcal M_{\rm NLO}
 =\mathcal M_{\rm vac.pol.}^{R}
 +\mathcal M_{{\rm vertex},e}^{R}
 +\mathcal M_{{\rm vertex},\mu}^{R}
 +\mathcal M_{\rm box}
 +\mathcal M_{\rm crossed}.
\]
This grouping is gauge invariant only after all members required at this
order are included.

\SupplementarySolution{2}{The General Feynman-Parameter Simplex and Massless Bubble in \(d\) Dimensions}
\begin{enumerate}
\item
For \(\Re\nu_i>0\),
\[
 A_i^{-\nu_i}=\frac{1}{\Gamma(\nu_i)}
 \int_0^\infty dt_i\,t_i^{\nu_i-1}e^{-t_iA_i}.
\]
Set \(t_i=T x_i\), where \(T\geq0\), \(x_i\geq0\), and
\(\sum_i x_i=1\).  The Jacobian is \(T^{n-1}\), so if
\(\nu=\sum_i\nu_i\),
\[
\begin{aligned}
 \prod_iA_i^{-\nu_i}
 ={}&\frac{1}{\prod_i\Gamma(\nu_i)}
 \int_0^\infty dT\,T^{\nu-1}\\
 &{}\times\int_0^1\prod_i dx_i\,
 \delta\left(1-\sum_i x_i\right)
 \prod_i x_i^{\nu_i-1}
 e^{-T\sum_i x_iA_i}.
\end{aligned}
\]
The \(T\)-integral is
\(\Gamma(\nu)(\sum_i x_iA_i)^{-\nu}\).  Therefore
\[
 \boxed{\;
 \frac{1}{\prod_iA_i^{\nu_i}}
 =\frac{\Gamma(\nu)}{\prod_i\Gamma(\nu_i)}
 \int_0^1\prod_i dx_i\,
 \frac{\delta(1-\sum_i x_i)\prod_i x_i^{\nu_i-1}}
 {(\sum_i x_iA_i)^\nu}.
 \;}
\]
For all \(\nu_i=1\), the normalization becomes \((n-1)!\), as required.
\item
The two-denominator identity gives
\[
 \frac1{A^aB^b}
 =\frac{\Gamma(a+b)}{\Gamma(a)\Gamma(b)}
 \int_0^1dx\,\frac{x^{a-1}(1-x)^{b-1}}
 {[xA+(1-x)B]^{a+b}}.
\]
After shifting the loop momentum and Wick rotating, the radial integral and
the remaining beta integral give
\[
 \boxed{\;
 I_B(a,b;p^2)
 =\frac{i}{(4\pi)^{d/2}}
 (-p^2-i0)^{d/2-a-b}
 \frac{\Gamma(a+b-d/2)\Gamma(d/2-a)\Gamma(d/2-b)}
 {\Gamma(a)\Gamma(b)\Gamma(d-a-b)} .
 \;}
\]
The phase follows from the stated \(i0\) prescription.  With
\(d=4-2\varepsilon\), \(a=b=1\), and a scale \(\mu\) inserted to keep the
integral dimensionless,
\[
 I_B(1,1;p^2)=\frac{i}{(4\pi)^2}
 \left[
 \frac1\varepsilon-\gamma+\ln4\pi
 +2-\ln\frac{-p^2-i0}{\mu^2}+O(\varepsilon)
 \right].
\]
The \(1/\varepsilon\) term is the logarithmic ultraviolet divergence.
\end{enumerate}

\SupplementarySolution{3}{Rank-Four Rotational Averaging and Tensor Reduction of a Bubble}
\begin{enumerate}
\item
An odd tensor integral vanishes because the measure and \(f(k^2)\) are even:
\[
 \int d^dk\,k^\mu k^\nu k^\rho f(k^2)=0.
\]
Rotational invariance and index symmetry imply
\[
 \int d^dk\,k^\mu k^\nu k^\rho k^\sigma f(k^2)
 =A(\delta^{\mu\nu}\delta^{\rho\sigma}
 +\delta^{\mu\rho}\delta^{\nu\sigma}
 +\delta^{\mu\sigma}\delta^{\nu\rho}).
\]
Contracting \(\mu\) with \(\nu\) and \(\rho\) with \(\sigma\) gives
\[
 \int d^dk\,(k^2)^2f(k^2)=A(d^2+2d),
\]
so
\[
 \boxed{\;
 \int d^dk\,k^\mu k^\nu k^\rho k^\sigma f(k^2)
 =\frac{\delta^{\mu\nu}\delta^{\rho\sigma}
 +\delta^{\mu\rho}\delta^{\nu\sigma}
 +\delta^{\mu\sigma}\delta^{\nu\rho}}
 {d(d+2)}
 \int d^dk\,(k^2)^2f(k^2).
 \;}
\]
The analogous rank-two coefficient is \(\delta^{\mu\nu}/d\).
\item
Write \(I=I_B(1,1;p^2)\).  The change of variables \(k\mapsto-k-p\)
immediately gives
\[
 I^\mu\equiv\int\frac{d^dk}{(2\pi)^d}
 \frac{k^\mu}{k^2(k+p)^2}=-\frac{p^\mu}{2}I.
\]
For
\(I^{\mu\nu}=A\eta^{\mu\nu}+Bp^\mu p^\nu\), contraction with the metric
leaves a scaleless tadpole:
\[
 dA+p^2B=\int\frac{d^dk}{(2\pi)^d}\frac1{(k+p)^2}=0.
\]
Also \(2p\cdot k=(k+p)^2-k^2-p^2\).  Squaring this relation inside the
integral shows, after scaleless terms are dropped, that
\[
 p_\mu p_\nu I^{\mu\nu}=\frac{p^4}{4}I.
\]
Solving the two linear equations gives
\[
 \boxed{\;
 I^{\mu\nu}
 =\frac{I}{4(d-1)}
 \left(dp^\mu p^\nu-p^2\eta^{\mu\nu}\right).
 \;}
\]
Its trace vanishes, consistently with dimensional regularization of the
scaleless tadpole.
\end{enumerate}

\SupplementarySolution{4}{Massless QED Self-Energy Poles}
In Feynman gauge, gamma algebra reduces the electron numerator to
\(\gamma_\mu(\sl p-\sl k)\gamma^\mu=(2-d)(\sl p-\sl k)\).
Feynman parametrization and the rank-one reduction then yield the pole
\[
 \Sigma_{\rm loop}^{\rm div}(p)
 =-\frac{e^2}{16\pi^2\varepsilon}\,i\sl p
\]
in the inverse-propagator convention used here.  It is cancelled by
\[
 Z_2=1-\frac{e^2}{16\pi^2\varepsilon}+O(e^4)
\]
with the associated sign determined by the counterterm placement in
\(S_R^{-1}\).

For the photon, the trace and symmetric integration give
\[
 \Pi_{\rm loop}^{\mu\nu}(q)\big|_{\rm div}
 =\frac{e^2}{12\pi^2\varepsilon}
 (q^2\eta^{\mu\nu}-q^\mu q^\nu).
\]
Thus \(q_\mu\Pi_{\rm loop}^{\mu\nu}=0\) before subtraction, and
\[
 Z_3=1-\frac{e^2}{12\pi^2\varepsilon}+O(e^4)
\]
cancels the pole.  The absence of a term proportional to
\(\eta^{\mu\nu}\) alone is the regulator-preserved statement that no photon
mass is induced.

\SupplementarySolution{5}{Euclidean Master Integrals and Shifted Tensor Integrals}
\begin{enumerate}
\item
Cartesian factorization gives
\[
 \int_{\mathbb R^d}d^dx\,e^{-x^2}
 =\left(\int_{-\infty}^{\infty}e^{-u^2}du\right)^d
 =\pi^{d/2}.
\]
In radial coordinates the same integral is
\[
 \Omega_d\int_0^\infty dr\,r^{d-1}e^{-r^2}
 =\frac{\Omega_d}{2}\Gamma(d/2).
\]
Equating them gives
\[
 \boxed{\qquad
 \Omega_d=\frac{2\pi^{d/2}}{\Gamma(d/2)}.
 \qquad}
\]
The right side is meromorphic in complex \(d\) and defines the continuation
used by dimensional regularization.  For positive integer \(d\), it agrees
with the geometric area; for general \(d\), only the analytic formula is
needed.
\item
Using radial coordinates and \(u=q^2/m^2\),
\[
\begin{aligned}
 J_n(m^2)
 &=\frac{\Omega_d}{(2\pi)^d}\int_0^\infty
 dq\,\frac{q^{d-1}}{(q^2+m^2)^n}\\
 &=\frac{\Omega_d}{2(2\pi)^d}(m^2)^{d/2-n}
 \int_0^\infty du\,\frac{u^{d/2-1}}{(1+u)^n}.
\end{aligned}
\]
The last integral is
\(B(d/2,n-d/2)\), so
\[
 \boxed{\;
 J_n(m^2)=\frac{1}{(4\pi)^{d/2}}
 \frac{\Gamma(n-d/2)}{\Gamma(n)}
 (m^2)^{d/2-n}.
 \;}
\]
Before continuation it converges for
\(\Re d>0\) and \(\Re n>\Re d/2\).  Ultraviolet poles occur when
\(n-d/2=0,-1,-2,\ldots\), exactly the poles of the Gamma function.
\item
Complete the square:
\[
 q^2+2q\cdot p+m^2=(q+p)^2+M^2,\qquad M^2=m^2-p^2.
\]
For \(M^2>0\) in the Euclidean convergence domain, translation invariance of
the regulated integral justifies \(\ell=q+p\).  Exercise S.11.8 then gives
\[
 \boxed{\;
 \int\frac{d^dq}{(2\pi)^d}
 \frac{1}{(q^2+m^2+2q\cdot p)^n}
 =\frac{\Gamma(n-d/2)}
 {(4\pi)^{d/2}\Gamma(n)}
 (m^2-p^2)^{d/2-n}.
 \;}
\]
The expression elsewhere is defined by analytic continuation with the
appropriate pole prescription.  A cutoff that obstructed the momentum shift
would not preserve this simple identity.
\item
After \(\ell=q+p\), odd powers of \(\ell\) integrate to zero.  Denoting the
scalar answer of S.11.9 by \(J_n(M^2)\), \(M^2=m^2-p^2\), we find
\[
 \boxed{\;
 \int\frac{d^dq}{(2\pi)^d}
 \frac{q^\mu}{(q^2+m^2+2q\cdot p)^n}
 =-p^\mu J_n(M^2).
 \;}
\]
For rank two,
\[
\begin{aligned}
 \int\frac{d^dq}{(2\pi)^d}
 \frac{q^\mu q^\nu}{(q^2+m^2+2q\cdot p)^n}
 ={}&p^\mu p^\nu J_n(M^2)\\
 &+\frac{\delta^{\mu\nu}}{d}
 [J_{n-1}(M^2)-M^2J_n(M^2)].
\end{aligned}
\]
Equivalently, the second line is
\[
 \frac{\delta^{\mu\nu}}{2(4\pi)^{d/2}}
 \frac{\Gamma(n-1-d/2)}{\Gamma(n)}
 (M^2)^{d/2+1-n}.
\]
Differentiating the scalar integral with respect to \(p_\mu\) produces the
same vector and tensor relations, providing an independent check.
\end{enumerate}

\SupplementarySolution{6}{The Spinor Vacuum-Polarization Numerator and Subtracted Vacuum Polarization at Small Momentum}
\begin{enumerate}
\item
The closed loop gives
\[
 \Pi^{\mu\nu}(p)=-e^2\int\frac{d^dk}{(2\pi)^d}
 \frac{\operatorname{Tr}[
 \gamma^\mu(-i\sl k+m)\gamma^\nu(-i(\sl k+\sl p)+m)]}
 {(k^2+m^2)((k+p)^2+m^2)}.
\]
The trace identity reduces the numerator to a sum of
\(k^\mu k^\nu\), \(k^{(\mu}p^{\nu)}\), and metric terms.  Introducing
\(x\), shifting \(\ell=k+xp\), and dropping odd powers gives
\[
 \Delta_x=m^2+x(1-x)p^2,
\]
\[
 \begin{aligned}
 \Pi^{\mu\nu}(p)
  &=-4e^2\int_0^1dx\int\frac{d^d\ell}{(2\pi)^d}
   \frac{N^{\mu\nu}(\ell,p;x)}{(\ell^2+\Delta_x)^2},\\
 N^{\mu\nu}(\ell,p;x)
  &=2\ell^\mu\ell^\nu-\eta^{\mu\nu}\ell^2
    -2x(1-x)p^\mu p^\nu
    +\eta^{\mu\nu}\Delta_x .
 \end{aligned}
\]
up to the common Wick-rotation factor fixed by convention.  Now use
\(\ell^\mu\ell^\nu\to\eta^{\mu\nu}\ell^2/d\).  The scalar master-integral
recursion cancels the remaining non-transverse metric term, leaving
\[
 \Pi^{\mu\nu}(p)
 =(p^2\eta^{\mu\nu}-p^\mu p^\nu)\Pi(p^2).
\]
The cancellation would be spoiled by an uncorrelated hard momentum cutoff.
\item
The on-shell photon normalization subtracts the unrenormalized scalar
coefficient at \(q^2=0\).  The pole and all scheme-dependent constants then
cancel, leaving
\[
 \boxed{\;
 \pi(q^2)=\frac{e^2}{2\pi^2}\int_0^1dx\,x(1-x)
 \ln\left(1+\frac{q^2x(1-x)}{m^2}\right).
 \;}
\]
For small \(q^2\), expand the logarithm.  The needed beta integrals are
\[
 \int_0^1x^2(1-x)^2dx=\frac1{30},\qquad
 \int_0^1x^3(1-x)^3dx=\frac1{140}.
\]
Therefore
\[
 \boxed{\;
 \pi(q^2)=\frac{e^2}{60\pi^2}\frac{q^2}{m^2}
 -\frac{e^2}{560\pi^2}\frac{q^4}{m^4}
 +O(q^6/m^6).
 \;}
\]
In particular \(\pi(0)=0\), and a heavy charged particle decouples from
low-momentum propagation as \(q^2/m^2\).
\end{enumerate}

\SupplementarySolution{7}{QED Counterterms and On-Shell Renormalization}
\begin{enumerate}
\item
The fields have dimensions
\([A_\mu]=1\), \([\psi]=3/2\), and \([\partial_\mu]=1\).  The local,
Lorentz- and gauge-invariant, \(C\)- and \(P\)-even operators of dimension at
most four that are already present in QED are
\[
 -\frac{\delta_3}{4}F_{\mu\nu}F^{\mu\nu},\qquad
 -\delta_2\bar\psi\gamma^\mu\partial_\mu\psi,\qquad
 \delta_m m\bar\psi\psi,\qquad
 -i\delta_1eA_\mu\bar\psi\gamma^\mu\psi.
\]
Gauge fixing renormalizes with the photon field and introduces no new
physical coupling.  A photon mass \(A_\mu A^\mu\) violates gauge
invariance.  A photon tadpole is odd under charge conjugation.  The Pauli
operator \(\bar\psi\sigma^{\mu\nu}F_{\mu\nu}\psi\) has dimension five, so
its finite coefficient can be predicted in renormalizable QED but it cannot
absorb a dimension-four ultraviolet divergence.  This census explains why
the counterterms in the chapter suffice.
\item
Current conservation makes the self-energy transverse:
\[
 \Pi^{\mu\nu}(q)
 =(q^2\eta^{\mu\nu}-q^\mu q^\nu)\pi(q^2).
\]
On the transverse subspace, Dyson summation is scalar and gives
\[
 D_T'(q)=\frac{D_T(q)}{1-\pi(q^2)}
 \propto\frac{1}{q^2[1-\pi(q^2)]}.
\]
Hence the massless pole has residue multiplied by
\([1-\pi(0)]^{-1}\).  The on-shell field normalization demands unit
residue, so
\[
 \boxed{\pi(0)=0.}
\]
At very small momentum transfer, the Coulomb amplitude is then
\(-e^2/q^2\) with \(e\) equal to the measured long-distance charge.  A
different finite subtraction would merely define a different parameter
\(e(\mu)\); converting it to the Thomson-limit charge restores the same
observable amplitude.
\item
Let \(z=i\sl p\).  Parity permits
\(\Sigma(p)=A(p^2)z+B(p^2)m\), and the full inverse propagator is
\[
 S_R^{-1}(p)=z+m-\Sigma_R(z).
\]
A pole at the measured mass requires
\[
 \boxed{\Sigma_R(-m)=0,}
\]
while unit residue requires
\[
 \boxed{\left.\frac{\partial\Sigma_R}{\partial z}\right|_{z=-m}=0.}
\]
The first condition fixes the mass counterterm and the second fixes the
field-strength counterterm.  Taylor expansion then gives
\[
 \Sigma_R(z)=O((z+m)^2),\qquad
 S_R^{-1}(p)=(z+m)[1+O(z+m)].
\]
Thus \(S_R(p)\) has a simple pole with residue one.  The phrase ``double
zero'' applies to the subtracted self-energy, not to the full inverse
propagator.
\end{enumerate}

\SupplementarySolution{8}{The Equality of Vertex and Wave-Function Poles}
The Ward--Takahashi identity for the proper vertex is
\[
 q_\mu\Gamma^\mu(p+q,p)
 =S^{-1}(p+q)-S^{-1}(p).
\]
Taking \(q\to0\) gives
\[
 \Gamma^\mu(p,p)=\frac{\partial S^{-1}(p)}{\partial p_\mu}.
\]
The local ultraviolet divergence on the left is proportional to
\(\delta_1\gamma^\mu\); the divergence on the right is the derivative of
the kinetic counterterm and is proportional to \(\delta_2\gamma^\mu\).
Therefore
\[
 \boxed{Z_1=Z_2}
\]
to one loop and, with a gauge-preserving regulator, to all orders.  Since
\(e_0=\mu^\varepsilon Z_1Z_2^{-1}Z_3^{-1/2}e\), the equality leaves
\[
 e_0=\mu^\varepsilon Z_3^{-1/2}e.
\]
Thus charge renormalization is determined entirely by photon
wave-function renormalization.

\SupplementarySolution{9}{Why Four-Photon Scattering Is Ultraviolet Finite}
Power counting assigns logarithmic superficial degree to a fermion loop with
four external photons.  Any remaining local divergence would therefore have
to be cancelled by a gauge-invariant operator of dimension four containing
four photon fields.  No such operator exists: gauge invariance requires
photons to occur through \(F_{\mu\nu}\), and an \(F^4\) operator has
dimension eight.  Consequently the logarithmic pieces of the six oriented
box orderings cancel, and the complete amplitude is ultraviolet finite.

For an odd number of photons, reversing the fermion-loop orientation gives
a factor \((-1)^n\) under charge conjugation.  At \(n=3\) the two
orientations cancel.  This is Furry's theorem.  The first nonzero
low-energy local interaction is therefore \(F^4/m^4\), consistent with both
claims.

\SupplementarySolution{10}{From the Pauli Form Factor to \(g-2\)}
For on-shell spinors, Lorentz covariance and current conservation give
\[
 \bar u(p')\Gamma^\mu u(p)
 =e\,\bar u(p')\left[
 \gamma^\mu F_1(q^2)
 +\frac{i\sigma^{\mu\nu}q_\nu}{2m}F_2(q^2)
 \right]u(p),
 \qquad q=p'-p.
\]
The Thomson-limit definition of charge fixes \(F_1(0)=1\).  In a slowly
varying magnetic field, the nonrelativistic Hamiltonian obtained from this
vertex contains
\[
 H_{\rm mag}=-\frac{e}{m}[F_1(0)+F_2(0)]
 \mathbf S\cdot\mathbf B.
\]
Comparing with
\(H_{\rm mag}=-g_e(e/2m)\mathbf S\cdot\mathbf B\) yields
\[
 g_e=2[1+F_2(0)],\qquad
 \boxed{\frac{g_e-2}{2}=F_2(0).}
\]
The Dirac form factor normalizes the charge; the Pauli form factor measures
the anomalous magnetic moment.

\SupplementarySolution{11}{Running Charge, Screening, and Heavy-Threshold Decoupling}
\begin{enumerate}
\item
The Ward identity gives
\[
 e_0=\mu^\varepsilon Z_3^{-1/2}e,\qquad d=4-2\varepsilon.
\]
For one Dirac fermion, minimal subtraction yields
\[
 Z_3=1-\frac{e^2}{12\pi^2\varepsilon}+O(e^4).
\]
Differentiate \(\ln e_0\) with respect to \(\ln\mu\) at fixed \(e_0\):
\[
 0=\varepsilon+\frac{\beta(e)}e
 -\frac12\beta(e)\frac{\partial\ln Z_3}{\partial e}.
\]
Solving perturbatively, while retaining the classical term until it
multiplies the pole, gives
\[
 \boxed{\beta(e)=-\varepsilon e+\frac{e^3}{12\pi^2}+O(e^5).}
\]
At \(d=4\), \(\varepsilon=0\), so the quantum term is positive.
\item
At one loop,
\[
 \frac{de}{d\ln\mu}=\frac{e^3}{12\pi^2}.
\]
Therefore
\[
 \frac{d}{d\ln\mu}\frac1{e^2}=-\frac1{6\pi^2},
\]
and integration gives
\[
 \boxed{\;
 \frac1{e^2(\mu_2)}
 =\frac1{e^2(\mu_1)}
 -\frac1{6\pi^2}\ln\frac{\mu_2}{\mu_1}.
 \;}
\]
Thus \(e(\mu)\) grows as the distance scale \(1/\mu\) decreases.  In
position space, virtual electron--positron pairs align so that charge of
opposite sign lies closer to a test charge.  The cloud screens the source at
long distance, making the infrared charge smaller than the short-distance
charge.  The positive beta function is the momentum-space expression of
this screening.
\item
For a fermion of mass \(M\), on-shell subtraction gives
\[
 \pi_R(q^2)=\frac{e^2}{2\pi^2}\int_0^1dx\,x(1-x)
 \ln\left(1+\frac{q^2x(1-x)}{M^2}\right).
\]
At low momentum,
\[
 \boxed{\;
 \pi_R(q^2)=\frac{e^2}{60\pi^2}\frac{q^2}{M^2}
 +O(q^4/M^4).
 \;}
\]
The heavy fermion therefore changes low-energy amplitudes only through
higher-dimension operators suppressed by \(M^{-2}\).

Minimal subtraction removes only poles and is insensitive to masses, so its
beta function formally includes the heavy fermion even for \(\mu\ll M\).
The effective theory resolves this by matching at \(\mu\simeq M\): a finite
threshold correction relates the couplings above and below \(M\), after
which the heavy field is omitted from the low-energy beta function.
\end{enumerate}

\SupplementarySolution{12}{The Euler--Heisenberg Effective Interaction}
\begin{enumerate}
\item
In four dimensions, the two independent parity-even scalars quartic in an
abelian field strength may be chosen as
\[
 \mathcal O_1=(F_{\mu\nu}F^{\mu\nu})^2,\qquad
 \mathcal O_2=(F_{\mu\nu}\widetilde F^{\mu\nu})^2,
\]
where
\(\widetilde F^{\mu\nu}=\tfrac12\epsilon^{\mu\nu\rho\sigma}F_{\rho\sigma}\).
Other contractions of four antisymmetric tensors reduce to these in four
dimensions.

With the mostly-plus convention,
\[
 F_{\mu\nu}F^{\mu\nu}=2(\mathbf B^2-\mathbf E^2),\qquad
 F_{\mu\nu}\widetilde F^{\mu\nu}=-4\mathbf E\cdot\mathbf B
\]
up to the sign fixed by the displayed dual convention; its square is
unambiguous.  Thus the two physical invariants are
\[
 (\mathbf E^2-\mathbf B^2)^2,\qquad
 (\mathbf E\cdot\mathbf B)^2.
\]
\item
For constant fields, introduce invariants \(a,b\) such that the eigenvalues
of \(F^\mu{}_\nu\) are encoded by an electric-like and magnetic-like pair.
After subtracting the zero-field vacuum term and the quadratic charge
renormalization term, the spinor determinant has the schematic form
\[
 \mathcal L_1=-\frac1{8\pi^2}\int_0^\infty\frac{ds}{s^3}e^{-m_e^2s}
 \left[(es)^2ab\,\coth(eas)\cot(ebs)-1
 -\frac{e^2s^2}{3}(a^2-b^2)\right].
\]
Using
\(\coth z=z^{-1}+z/3-z^3/45+\cdots\) and
\(\cot z=z^{-1}-z/3-z^3/45+\cdots\), the first surviving term is quartic.
Since \(\int_0^\infty ds\,s\,e^{-m_e^2s}=m_e^{-4}\), one obtains
\[
 \boxed{\;
 \mathcal L_{\rm EH}
 =\frac{2\alpha^2}{45m_e^4}
 \left[(\mathbf E^2-\mathbf B^2)^2
 +7(\mathbf E\cdot\mathbf B)^2\right].
 \;}
\]
Equivalently,
\[
 \mathcal L_{\rm EH}
 =\frac{\alpha^2}{90m_e^4}
 \left[(F_{\mu\nu}F^{\mu\nu})^2
 +\frac74(F_{\mu\nu}\widetilde F^{\mu\nu})^2\right].
\]
\item
Each external photon enters \(\mathcal L_{\rm EH}\) through one field
strength, so it contributes one power of its momentum.  The coefficient of
the four-photon vertex is \(O(\alpha^2/m_e^4)\).  Consequently
\[
 \boxed{\mathcal M(\gamma\gamma\to\gamma\gamma)
 =O\left(\alpha^2\frac{\omega^4}{m_e^4}\right).}
\]
For relativistic \(2\to2\) scattering, the cross section scales as
\(\lvert\mathcal M\rvert^2/s\), with \(s=O(\omega^2)\).  Hence
\[
 \boxed{\sigma=O\left(
 \alpha^4\frac{\omega^6}{m_e^8}\right).}
\]
No lower power is allowed: gauge invariance requires every external photon
to appear through \(F_{\mu\nu}\), while locality below threshold makes the
amplitude analytic in momenta.  The next corrections contain additional
derivatives and are further suppressed by \(\omega^2/m_e^2\).
\end{enumerate}

\SupplementarySolution{13}{Pauli Operators, Nonrelativistic Reduction, and Simplex Parameters}
\begin{enumerate}
\item
Define
\(\sigma^{\mu\nu}=\tfrac{i}{2}[\gamma^\mu,\gamma^\nu]\).  The on-shell
vertex at small momentum transfer contains
\[
 e\Gamma^\mu=e\gamma^\mu
 +e\frac{i\sigma^{\mu\nu}q_\nu}{2m}F_2(0)+O(q^2).
\]
The local interaction
\[
 \boxed{\;
 \Delta\mathcal L
 =\frac{eF_2(0)}{4m}
 \bar\psi\sigma^{\mu\nu}F_{\mu\nu}\psi
 \;}
\]
produces precisely this vertex: differentiating \(F_{\mu\nu}\) with
respect to \(A_\mu\) supplies two equal antisymmetric contributions and one
factor of \(iq_\nu\).  The operator has dimension five, so its coefficient
has dimension \(-1\).  It is gauge invariant and \(C\)-, \(P\)-, and
\(T\)-even for a magnetic dipole; replacing it by the dual field strength
would instead describe an electric dipole and violate \(P\) and \(T\).
\item
Write the Dirac spinor after removing its rest-energy phase as
\(\psi=e^{-imt}(\Phi,X)^{\mathsf T}\), and put
\(\boldsymbol\pi=\mathbf p-e\mathbf A\).  To leading order the lower
component obeys
\[
 X=\frac{\boldsymbol\sigma\cdot\boldsymbol\pi}{2m}\Phi+O(m^{-2}).
\]
Substitution into the upper-component equation uses
\[
 (\boldsymbol\sigma\cdot\boldsymbol\pi)^2
 =\boldsymbol\pi^2+e\boldsymbol\sigma\cdot\mathbf B
\]
with the sign adjusted if \(e\) denotes the positive magnitude of the
electron charge.  Including the Pauli operator gives
\[
 i\partial_t\Phi=
 \left[
 \frac{\boldsymbol\pi^2}{2m}+eA^0
 +\frac{e}{2m}[1+F_2(0)]
 \boldsymbol\sigma\cdot\mathbf B
 \right]\Phi.
\]
Thus, in terms of \(\mathbf S=\boldsymbol\sigma/2\),
\[
 H_{\rm mag}=-g_e\frac{e}{2m}\mathbf S\cdot\mathbf B,\qquad
 \boxed{g_e=2[1+F_2(0)]},
\]
where the overall sign follows the particle's signed charge.
\item
For denominators \(A_1,A_2,A_3\),
\[
 \frac1{A_1A_2A_3}
 =2\int_0^1\prod_{i=1}^3dx_i\,
 \frac{\delta(1-x_1-x_2-x_3)}
 {(x_1A_1+x_2A_2+x_3A_3)^3}.
\]
Take \(A_1=(k+p')^2+m^2-i0\),
\(A_2=(k+p)^2+m^2-i0\), and \(A_3=k^2-i0\), with
\(p^2=p'{}^2=-m^2\) and \(q=p'-p\).  Shifting
\(\ell=k+x_1p'+x_2p\) gives
\[
 \boxed{\;
 x_1A_1+x_2A_2+x_3A_3
 =\ell^2+m^2(x_1+x_2)^2+q^2x_1x_2-i0.
 \;}
\]
Since \(x_1+x_2=1-x_3\), this is
\(\ell^2+m^2(1-x_3)^2+q^2x_1x_2-i0\).  The common form makes the finite
Pauli projection manifest.
\end{enumerate}

\SupplementarySolution{14}{The One-Loop Pauli Form Factor}
\begin{enumerate}
\item
After the shift in S.11.13(c), terms odd in \(\ell\) vanish.  Symmetric
integration replaces
\(\ell^\alpha\ell^\beta\) by
\(\eta^{\alpha\beta}\ell^2/d\).  Gamma identities reduce the even
\(\ell^2\) term to a multiple of \(\gamma^\mu\), so it contributes only to
the Dirac form factor and its ultraviolet subtraction.

For the remaining terms, move \(\sl p\) and \(\sl p'\) onto the external
spinors and use
\((i\sl p+m)u(p)=0\) and
\(\bar u(p')(i\sl p'+m)=0\).  The part independent of \(\ell\) that is not
proportional to \(\gamma^\mu\) reduces to
\[
 \boxed{\;
 N^\mu_{\rm Pauli}
 =2im(x_1+x_2)x_3(p^\mu+p'{}^\mu).
 \;}
\]
The Gordon identity rewrites \(p^\mu+p'{}^\mu\) between on-shell spinors as
a combination of \(\gamma^\mu\) and
\(i\sigma^{\mu\nu}q_\nu\).  After the \(F_1(0)\) piece is fixed by charge
renormalization, the remainder is exactly the Pauli form factor.
\item
Radial integration in four Euclidean dimensions gives
\[
\begin{aligned}
 \int\frac{d^4k}{(2\pi)^4}\frac1{(k^2+\Delta)^3}
 &=\frac{2\pi^2}{(2\pi)^4}
 \int_0^\infty dk\,\frac{k^3}{(k^2+\Delta)^3}\\
 &=\boxed{\frac1{32\pi^2\Delta}}.
\end{aligned}
\]
Combining this with the numerator from S.11.14(a) yields the finite simplex
representation
\[
 \boxed{\;
 F_2(q^2)=\frac{e^2}{4\pi^2}
 \int_0^1\prod_{i=1}^3dx_i\,
 \delta(1-\textstyle\sum_i x_i)
 \frac{m^2x_3(1-x_3)}
 {m^2(1-x_3)^2+q^2x_1x_2}.
 \;}
\]
No ultraviolet regulator remains: the Pauli projection has lowered the
superficial degree sufficiently for convergence.
\item
At \(q^2=0\), the integrand in S.11.14(b) becomes
\[
 \frac{x_3}{1-x_3}.
\]
For fixed \(x_3\), the line segment
\(x_1+x_2=1-x_3\) has parameter measure \(1-x_3\).  Therefore
\[
\begin{aligned}
 F_2(0)
 &=\frac{e^2}{4\pi^2}\int_0^1dx_3\,
 (1-x_3)\frac{x_3}{1-x_3}\\
 &=\frac{e^2}{4\pi^2}\int_0^1x_3\,dx_3
 =\frac{e^2}{8\pi^2}
 =\boxed{\frac{\alpha}{2\pi}}.
\end{aligned}
\]
Using \(g_e=2[1+F_2(0)]\) then gives
\[
 \boxed{\qquad g_e=2+\frac{\alpha}{\pi}+O(\alpha^2),\qquad
 \frac{g_e-2}{2}=\frac{\alpha}{2\pi}+O(\alpha^2).\qquad}
\]
\end{enumerate}
\end{enumerate}
